\section{Counting common zeros by polynomial interpolation over \texorpdfstring{$\mathbb Q$}{Q}}\label{sec:setup}
We now prove Theorem~\ref{thm:overview-sparse-polynomials}. The input
is a system $\mathcal P$ of $m_{\mathrm{test}}\le cn$ polynomial equations
$P_i(x)=0$ over $\mathbb F_2$ in $n$ Boolean variables. Each $P_i$ is a
list of at most $(\log n)^c$ monomials of degree at most $k$, for fixed
$k,c$. Our target is $N(\mathcal P)$ from \eqref{eq:common-zero-count}.
Use the integer lifts $L_i$, scopes $I_i$, linear coefficient polynomials $e_j$,
and coefficient algebra $\mathcal E_\ell$ from
Sections~\ref{subsec:rational-lift}--\ref{subsec:total-order}, with
$\ell=n+1$. Lemma~\ref{lem:total-order} bounds the product degree by
$k(m_{\mathrm{test}}+2n)=O_{k,c}(n)$. On each branch fixing the few
variables of high incidence, Lemmas~\ref{lem:incidence-parameters}
and~\ref{lem:layout} supply the data used by Proposition~\ref{prop:count}.
The full deduction, including the cost of these branches, is in
Section~\ref{subsec:sparse-queries}.

\subsection{Counting with supplied replacement forms}\label{subsec:count-polynomial}
For equation $i$ supply Section~\ref{sec:counting} with the factor
\begin{equation}\label{eq:lifted-factor}
 \mathcal B_i((X_v)_{v\in I_i},V)
 =\sum_{j=0}^{\ell-1}e_j(L_i(\vec X))h(L_i(\vec X))^j V^j
 \quad\text{in }\mathcal E_\ell[(X_v)_{v\in I_i}].
\end{equation}
Within a factor, every occurrence of $X_v$ uses the same replacement
form; different factors may use different forms for $X_v$. Even if
each factor receives its own independent inputs, $L_i$ has degree at
most $k$ in those inputs. Lemma~\ref{lem:total-order} bounds the input
degree of the product \emph{in $\mathcal E_\ell$} by
$d=k(m_{\mathrm{test}}+2(\ell-1))$. Thus the product satisfies precisely
the degree hypothesis \eqref{eq:general-product-degree} of the common
routine. Adding factors equal to one with scope $\{v\}$ ensures that
all declared variables occur in the construction.

Let $H(T)\in\mathcal E_\ell[T]$ be \eqref{eq:general-count-polynomial}
for these factors. Its $V$-coefficients at $T=0$ are integers because
all replacements restore their Boolean inputs at zero. By the same
lemma, summing their values gives
\begin{equation}\label{eq:origin-residue}
 \deg_T H\le d_H:=kR(m_{\mathrm{test}}+2(\ell-1)),\qquad
 \mathcal S_\ell(H(0))\equiv N(\mathcal P)\pmod{2^\ell}.
\end{equation}
Consequently the least nonnegative residue of $\mathcal S_\ell(H(0))$
is the exact count. Interpolation takes place over the rationals;
the congruence is needed only on Boolean inputs at $T=0$.

\subsection{Cost of evaluating the factors}\label{subsec:one-query-cost}
\begin{proposition}\label{prop:count}
Let $\mathcal P$ be a list of $m_{\mathrm{test}}$ polynomial equations
over $\mathbb F_2$ on $n$ declared Boolean variables, each given by
a monomial list of degree at most $k$. Suppose the factors
\eqref{eq:lifted-factor}, including the neutral factors, are supplied
with the data of Definition~\ref{def:layout}. Let $S\ge1$ bound
the total length of the monomial lists and factor scopes, and let
$\ell$ be any positive integer with $2^\ell>2^n$. With $d_H$ from
\eqref{eq:origin-residue} and $q,J$ from \eqref{eq:hasse-order},
$N(\mathcal P)$ can be computed deterministically in
\[
 \poly(n,m_{\mathrm{test}},\ell,D,Q,R,J,r,k,S)\,
 \left(\sum_{s=0}^{\min(q,n)}2^s\binom ns\right)
 (2D+1)^{2\tau n}2^{2n/r}
\]
bit operations and $\poly(n,m_{\mathrm{test}},\ell,D,Q,R,J,r,k,S)$
bits of working space.
\end{proposition}
\begin{proof}
Apply Proposition~\ref{prop:general-count} with coefficients in
$\mathcal E_\ell$. At an evaluation point $t$, the routine uses
$\mathcal E_\ell[z]/(z^J)$, so a stored element has just $J\ell$
rational coefficients. To evaluate factor $i$ at a fixed Boolean
substitution $\vec X=\mathbf1_U$ in \eqref{eq:taylor-substitution}, first
evaluate each listed monomial by multiplying its at most $k$ substituted
inputs, then add the results to get $L_i$. Use this value in
\eqref{eq:lifted-factor}, successively multiplying by $h(L_i)$ to
obtain its powers. This takes polynomially many rational operations in
the displayed parameters and never enumerates assignments to the entire
scope $I_i$. The polynomials $e_j$ are computed once by
Lemma~\ref{lem:rational-idempotent}.

For bit complexity, the rows have common denominator
$d_{\mathrm{den}}=(R-1)!Q!$. Each substituted monomial has denominator
dividing $\Delta=d_{\mathrm{den}}^k$, every lifted factor has denominator
dividing $\Gamma=\Delta^{2\ell-1}$, and their product has denominator
dividing $\Omega=\Gamma^{m_{\mathrm{test}}}$. Thus exponentially many
summands can be accumulated over one denominator.
Appendix~\ref{app:arithmetic} bounds the bit lengths in these sums by
\begin{equation}\label{eq:sum-bit-bound}
 B_{\mathrm{sum}}=\poly(n,m_{\mathrm{test}},\ell,k,Q,\log(D+2),\log(S+2)).
\end{equation}
For each $0\le s<\ell$ we interpolate the rational polynomial
$[V^s]H(T)$ from its $QJ$ Taylor coefficients. The appendix verifies
polynomial bit cost for this interpolation and the final integer
remainder in \eqref{eq:origin-residue}. The $O(QJ\ell)$ stored
coefficients, the supplied rows, partitions, assignments, and measurement matrices, and the sequential enumerations in
Proposition~\ref{prop:general-count} give the stated space bound.
\end{proof}

\subsection{Constructing the forms and completing the proof}\label{subsec:sparse-queries}
Each sparse polynomial reads few variables, but one variable might
occur in many equations. We first fix the values of the few variables
with the largest numbers of occurrences. On each resulting branch we
can construct the rows, partitions, assignments, and measurement matrices required by Proposition~\ref{prop:count}.

\begin{proof}[Proof of Theorem~\ref{thm:overview-sparse-polynomials}]
Let $\mathcal P$ be the supplied polynomial system. An empty list of
equations gives $2^n$ common zeros; a constant equation $1=0$ gives
none. Discard equations $0=0$. Include a neutral factor equal to one
with scope $\{v\}$ for each declared variable $v$, and write
\[
 M=\sum_{f\in\mathcal F}|I_f|
\]
for the total number of factor--variable incidences. The input
bounds imply $M=n(\log n)^{O_{k,c}(1)}$, and every scope has
$(\log n)^{O_{k,c}(1)}$ variables. A particular variable, however,
may occur in many of these scopes.

Put $\lambda=\lceil\log_2 n\rceil$ and
\[
 b_0=\left\lceil\frac Mn\right\rceil\lceil\sqrt\lambda\rceil.
\]
Let $S\subseteq[n]$ be the variables belonging to more than $b_0$
scopes. For each $b\in\{0,1\}^S$, we fix $x|_S=b$ and count the
completions on $[n]\setminus S$; we call this restricted problem a
branch. Since $|S|\le M/b_0\le n/\sqrt\lambda$, there are at most
$2^{n/\sqrt\lambda}$ branches.
Each branch has the same number $n'$ of free variables, with
\[
 n'\ge n-\frac n{\sqrt\lambda}\ge\frac n2
 \qquad(n\ge20).
\]
Substitute the fixed bits in every $P_i$ over $\mathbb F_2$ and
combine duplicate monomials modulo two before forming its integer
lift. Neither degree nor the number of monomials increases. A constant
equation can again be handled immediately. The branches partition the
Boolean cube, so their common-zero counts add to $N(\mathcal P)$.

To apply Section~\ref{sec:construction}, remove the fixed variables from the
scopes and relabel the free variables by $[n']$.
Write $I'_f\subseteq[n']$ for the restricted
scopes and $M'=\sum_f|I'_f|$ for their total incidence count.
Let $b$ be the maximum factor incidence of a free variable
and let $w=\max_f|I'_f|$. Then $b\le b_0$,
and both $b$ and $w$ are bounded by a fixed power of $\log n'$.
The remaining neutral factors ensure $M'\ge n'$.
Set
\[
 \lambda'=\lceil\log_2 n'\rceil,\qquad
 \nu'=\lceil\log_2\lambda'\rceil,
\]
and choose the smallest integer $\beta\ge1$ with
$b,w\le(\lambda')^\beta$. For fixed $k,c$, the possible values
of $\beta$ are bounded by a constant for all sufficiently large
$n$. We apply Lemma~\ref{lem:incidence-parameters} to the restricted
scopes $(I'_f)$ on $[n']$, with incidence count $M'$, using
\[
 r=\left\lfloor\frac{\sqrt{\lambda'}}{(\nu')^2}\right\rfloor,
 \qquad Q=\left\lfloor\frac{\lambda'}{128\nu'}\right\rfloor,
 \qquad R=(\beta+3)\lceil\sqrt{\lambda'}\rceil.
\]
If $r<2$ or $Q\le R$, use the exhaustive-search fallback below
before forming $\tau$, $\epsilon$, or the parity interpolation
interval $\mathcal I$.
Otherwise set
\[
 \tau=\frac1{r\lambda'},\qquad
 \epsilon=\frac1{br^2\lambda'},\qquad D=r^{2Q},
\]
and carry out the remaining feasibility checks from that lemma.
For all sufficiently large $n$, Lemma~\ref{lem:layout} supplies
the required rows and measurement matrices, with
\[
 r=\Theta\!\left(\frac{\sqrt{\log n'}}{(\log\log n')^2}\right),
 \qquad Q=\Theta\!\left(\frac{\log n'}{\log\log n'}\right),
 \qquad R=O_{k,c}(\sqrt{\log n'}),
 \qquad D\le(2n')^{1/128}.
\]

On every branch we may use the original precision $\ell=n+1$,
since even the unrestricted count is smaller than $2^\ell$. As
$n'\ge n/2$, both $\ell$ and $m_{\mathrm{test}}$ are $O_c(n')$.
Thus $d_H=O_{k,c}(n'R)$, $J=O_{k,c}(n'R/Q)+1$, and
\eqref{eq:three-costs}, in dimension $n'$, has exponent
\[
 O_{k,c}\!\left(\frac{n'(\log\log n')^2}{\sqrt{\log n'}}\right)
 =O_{k,c}\!\left(\frac{n(\log\log n)^2}{\sqrt{\log n}}\right).
\]
All substitution data, coefficient arrays, and interpolation systems
have total representation length $n^{O_{k,c}(1)}$ bits.
Multiplying the per-branch running time by at most
$2^{n/\sqrt\lambda}$ leaves the displayed exponential bound
unchanged, so summing the branch answers gives the required running time.

To implement the construction for every input size, the algorithm
checks its numerical hypotheses using integer arithmetic.
In addition to the preliminary guards above, it checks
$2(b+w)D^2\le M'$ and $2Qr^2\le n'$ from
Lemma~\ref{lem:incidence-labels}, the endpoints of $\mathcal I$,
$R\ge32$, $6R\le Q$, and the binomial tail bound in
\eqref{eq:incidence-tail}, which can be tested exactly.
If any check fails, the algorithm abandons the branched computation
and checks the original equations directly on all $2^n$ assignments. For fixed
$k,c$, the bounds just proved and
Lemma~\ref{lem:incidence-parameters} show that a failure can occur
only for finitely many $n$. This fallback is therefore absorbed
in the theorem's constants.

To combine the branch counts within the stated space bound, we
process the branches one at a time. We initialize an integer
accumulator $Z=0$ and add each branch's answer to it. The construction and
Proposition~\ref{prop:count} use $n^{O_{k,c}(1)}$ bits of space on
each branch.
Disjointness gives $0\le Z\le2^n$ throughout, and the final value
is $N(\mathcal P)$, as claimed.
\end{proof}

